\textbf{Navigation benchmarks.} BARN \cite{perille2020barn} provides procedurally generated
environments ordered by difficulty metrics and evaluates planners in Gazebo. MRPB
\cite{wen2020mrpb} evaluates local planners in Gazebo on safety, efficiency and smoothness
metrics. Bench-MR \cite{heiden2021benchmr} benchmarks sampling-based planners for wheeled robots,
building on OMPL's benchmarking infrastructure \cite{sucan2012ompl}. Arena-Bench
\cite{kastner2022arenabench} targets highly dynamic environments and integrates with ROS. These
suites run where their software stack is installed; \bench{} instead makes the run itself
portable and checkable (hash chains across engines), and treats the full ROS~2 stack as a
second, slower evidence class to validate against rather than as the only one.

\textbf{Difficulty.} Like BARN, we order scenarios by measurable features; ours come from a
reference plan (path length, narrowest clearance, bends, zones crossed, exposure to moving
actors) and are used for reporting per tier, not as a learned difficulty score.

\textbf{Reproducibility.} Calls for replicable, measurable robotics research
\cite{bonsignorio2015replicable} motivate our design decisions: frozen and versioned suites,
committed reference results, statistics with intervals \cite{wilson1927}, and no number in this
paper that is not generated from committed data.

\textbf{Algorithms.} The built-in references are breadth-first search, Dijkstra, A*, greedy and
weighted A*, and D* Lite \cite{koenig2002dstarlite} for planning; DWA \cite{fox1997dwa}, Regulated
Pure Pursuit \cite{macenski2023rpp} and MPPI \cite{williams2017mppi} for local control; and Monte
Carlo localisation \cite{fox1999mcl} and an EKF for localisation. The community entry in our
case study is Lazy Theta* \cite{nash2010lazytheta}, an any-angle variant of Theta*
\cite{nash2007theta}.
